% =============================================================================
%  s03a-framing.tex —— 第 3 节「标架化与路径积分」：标题、引言、3.1、3.2
%  覆盖原文 p.13（"Apologia for geometry."）至 p.19 式 (3.24)；编号公式 (3.1)–(3.24)
% =============================================================================

% ==== 原文 p.13 ====
\section{标架化与路径积分}\label{sec:3}

% [Tutor Note: 本节重点] 物理图像/数学工具/主要结论
\begin{keytakeaways}
\begin{itemize}[leftmargin=1.4em,itemsep=2pt,topsep=2pt]
  \item \textbf{物理图像：}非阿贝尔陈--西蒙斯作用量在规范变换下只差 $2\pi k\,w(g)$；绕数非零的大规范变换迫使能级 $k$ 取整数，量子理论靠振幅 $e^{i\SCS}$ 的单值性存活。
  \item \textbf{数学工具：}主 $G$-丛上的联络与曲率 $F=\dd A+A\wedge A$、博特定理 $\pi_3(G)=\Z$、法捷耶夫--波波夫规范固定、一圈行列式 $\det(\Delta)/\sqrt{\det(L_-)}$，以及阿蒂亚--帕托迪--辛格指标定理给出的 $\eta$ 相位。
  \item \textbf{主要结论：}大 $k$ 下配分函数由平坦联络的稳相求和给出：$\mathcal{Z}=e^{i\pi\eta[0]/2}\sum_\alpha e^{i(k+c_2(G)/2)I[A^{(\alpha)}]}T_\alpha$，能级平移 $k\to k+c_2(G)/2$ 在此出现。
\end{itemize}
\end{keytakeaways}

\paragraph{为几何辩护}
% 原文标题：Apologia for geometry.
上一节我们刻意保持朴素，选择了初等、却未必深刻的表述方式。现在，某种意义上要一切从头再来。这一回的叙述会明显更形式化，但这里发展出的几何始终由阿贝尔理论中积累的直觉引导。此前遇到的全部现象，如今都要配上纤维、丛与表示，抽象地铺展开来。我们先完整地建立非阿贝尔陈--西蒙斯（CS）理论，给出它在规范变换下的行为，再用纽结描述它的物理可观测量（威尔逊圈）。随后转向大 $k$ 下配分函数的鞍点计算；此时理论弱耦合。为正规化过程中出现的那些相当微妙的发散，我们将讨论定向 3-流形与纽结的标架化。最后，弄清路径积分在标架变换下如何变化，就能完成这些计算。

\subsection{非阿贝尔陈--西蒙斯作用量}\label{sec:3.1}

\paragraph{几何设定}
% 原文标题：Geometrical setting.
故事始于一个代表时空的定向 3-流形 $M$，以及一个紧致单李群 $G$——规范群，其李代数为 $\mathfrak{g}$。作为物理学家，我们有时把李群与它的李代数混为一谈。这里我们虽努力区分 $G$ 与 $\mathfrak{g}$，却难免仍有不精确之处。为使规范变换的概念严格化，固定一个主 $G$-丛\footnote{为方便起见，可以设想 $E=M\times G$ 是平凡丛，取 $M=S^3$、$G=\mathrm{SU}(N)$。我们的形式适用于更一般的情形，但这是遇到困难时值得回头参考的有用例子。} $E\to M$。$E$ 的截面称为规范变换，它们是映射 $g\colon M\to E$，把每个 $x\in M$ 光滑地对应到一点 $g(x)=(x,g_x)\in E$，其中 $g_x\in G$。我们常滥用记号，写成 $g(x)=g_x$，把 $g$ 看作映射 $M\to G$，给出依赖时空位置的 $G$ 元素。当 $g_x$ 接近 $G$ 的单位元时，我们改而考虑无穷小（生成元形式的）规范变换，并把 $g(x)$ 视为 $\mathfrak{g}$ 值的零形式，即依赖时空位置的李代数生成元。

接下来，在 $E$ 上引入主联络。在物理学家看来，\footnote{关于恰当的数学表述（包括伴丛方面的洞见），参见 \cite{steenrod1951} 或 \cite{nakahara}。} 它不过是 $M$ 上一个取 $\mathfrak{g}$ 值的 1-形式。记该联络为 $A$；在 $M$ 的局部坐标下写 $A=A_\mu(x)\dd x^\mu$。每个 $A_\mu(x)$ 都落在李代数里，因此给定 $\mathfrak{g}$ 的一组基 $T^a$，还可分解 $A_\mu(x)=A^a_\mu(x)T^a$，其中规范指标 $a\in\{1,\dots,d=\dim G\}$。这个对象 $A_\mu(x)$ 称为规范场，它推广了 \S\ref{sec:2} 的 $\mathrm{U}(1)$ 规范场。现在可以问一个自然的问题：既然 $E$ 自带 $G$ 作用，

% ==== 原文 p.14 ====
联络在该作用下如何变化？换句话说，规范场在规范变换下如何改变？答案就是下述变换律：
\begin{equation}\tag{3.1}\label{eq:3.1}
  A_\mu \longrightarrow A'_\mu = g^{-1}A_\mu g + g^{-1}\partial_\mu g,
  \qquad g = g(x)\in E.
\end{equation}

\paragraph{协变导数与曲率}
% 原文标题：Covariant derivative and curvature.
为得到该变换律的无穷小版本，引入规范协变导数 $D$。它作用在 $\mathfrak{g}$ 值微分形式 $\omega$ 上，$D_\mu\omega=\partial_\mu\omega+[A_\mu,\omega]$，其中 $[\cdot,\cdot]$ 是李括号。在局部规范变换 $\varepsilon(x)\in\Omega^0(M,\mathfrak{g})$ 下，规范场按协变导数改变：$A_\mu\to A_\mu-D_\mu\varepsilon$。用行话说，联络通过 $[A_\mu,\cdot]$ 把普通导数 $\partial_\mu$“扭”了一下。要考察 $A_\mu$ 引入了多少扭曲，可以测量 $D_\mu$ 与自身不对易的程度；若规范丛平坦，$D_\mu$ 本应对易。因此，我们如下定义联络的曲率，也称场强：
\begin{equation}\tag{3.2}\label{eq:3.2}
  F_{\mu\nu} \equiv [D_\mu, D_\nu] = \partial_\mu A_\nu - \partial_\nu A_\mu + [A_\mu, A_\nu]
  \quad\Longleftrightarrow\quad F = \dd A + A\wedge A.
\end{equation}

% [Tutor Note: 易错点] (3.2) 分量式中 [A_μ,A_ν] 的系数是 1 不是 1/2；形式式 F=dA+A∧A 里那 1/2 已被 A∧A 的反对称化吸收，两套写法不可重复记账。
\begin{yizhu}[易错点]
\eqref{eq:3.2} 的分量式里 $[A_\mu,A_\nu]$ 的系数是 $1$，不是 $1/2$。用形式记号写就是 $F=\dd A+A\wedge A$；分量式里的 $[A_\mu,A_\nu]$ 与形式式里的 $A\wedge A$ 只差一次反对称化，那个 $1/2$ 已经并进 $A\wedge A$ 里。反过来，从形式式走到分量式 $F_{\mu\nu}=\partial_\mu A_\nu-\partial_\nu A_\mu+[A_\mu,A_\nu]$ 时，切不可再补一个 $1/2$。这个定义还与协变导数 $D_\mu\omega=\partial_\mu\omega+[A_\mu,\omega]$ 绑定：若改用厄米约定写 $A$，$[A_\mu,A_\nu]$ 前的符号也要随之改，否则 $[D_\mu,D_\nu]$ 展不出来。
\end{yizhu}

若 $G$ 阿贝尔（例如 $G=\mathrm{U}(1)$），则所有对易子为零。此时 $D_\mu=\partial_\mu$，$F_{\mu\nu}$ 化为通常的麦克斯韦场强张量。但非阿贝尔场强含有 $A$ 的二次项，一般是非线性的。若 $F=0$，则称联络 $A$ 平坦。受 \eqref{eq:2.2} 启发，我们很快会看到：$E$ 上平坦联络的空间决定纯陈--西蒙斯理论的经典相空间。

\paragraph{陈--西蒙斯作用量}
% 原文标题：The Chern--Simons action.
从 $A$ 出发，我们如下定义陈--西蒙斯作用量：
\begin{equation}\tag{3.3}\label{eq:3.3}
\begin{aligned}
  \SCS[A] &\equiv \frac{k}{4\pi}\int_M \Tr\Bigl(A\wedge\dd A + \frac{2}{3}A\wedge A\wedge A\Bigr)\\
  &= \frac{k}{4\pi}\int_M \dd^3x\,\varepsilon^{\mu\nu\rho}\Tr\Bigl(A_\mu(\partial_\nu A_\rho - \partial_\rho A_\nu) + \frac{2}{3}A_\mu[A_\nu,A_\rho]\Bigr).
\end{aligned}
\end{equation}

% [Tutor Note: 原文疑误] (3.3) 第二行比第一行恰大一倍：A∧dA=½εA(∂νAρ−∂ρAν)、A∧A∧A=½εA[Aν,Aρ]，故括号前应有整体 ½，系数应为 1/2 与 1/3；且与 (3.4) 的 1/3 自相矛盾。
\begin{yizhu}[原文疑误]
按形式式展开，$\dd^3x\,\varepsilon^{\mu\nu\rho}$ 前应带整体 $\frac12$：$A\wedge\dd A=\frac12\varepsilon^{\mu\nu\rho}A_\mu(\partial_\nu A_\rho-\partial_\rho A_\nu)$，$A\wedge A\wedge A=\frac12\varepsilon^{\mu\nu\rho}A_\mu[A_\nu,A_\rho]$。照此，\eqref{eq:3.3} 第二行括号内的系数应是 $\frac12$ 与 $\frac13$，原文却写成 $1$ 与 $\frac23$，整行恰为第一行的两倍。旁证是 \eqref{eq:3.4} 与 \eqref{eq:3.5}：那里 $\frac13$ 与 $1/24\pi^2$ 配套得出 $\Delta\SCS=2\pi k\,w(g)$，只有按第一行解释才自洽。中译本照录原文，不作改动。
\end{yizhu}

$\mathfrak{g}$ 值的 3-形式 $\omega_{\mathrm{CS}}\equiv A\wedge\dd A+\frac23 A\wedge A\wedge A$ 就是著名的陈--西蒙斯形式；这里的 $\Tr$ 指 $\mathfrak{g}$ 上基灵型的某个倍数（归一化稍后确定）。更具体地说，在 $\mathfrak{g}$ 的伴随表示中，$\omega_{\mathrm{CS}}$ 不过是一个（由 3-形式组成的）矩阵，而这个矩阵的迹正比于 $\Tr(\omega)$。由 \eqref{eq:3.3} 得出的运动方程，来自场变分 $\delta A$ 所引起的 $\delta\SCS$，即 $\delta\SCS=0\Longrightarrow F^a_{\mu\nu}=0$。正如许诺的那样，这描述了 $E$ 上平坦联络的空间。

% [Tutor Note: 易错点] δS_CS=(k/2π)∫Tr(δA∧F) 只是体项；M 带边界时还有 (k/4π)∫_{∂M}Tr(A∧δA)，作者写 δS_CS=0\$\\ 即默认边界项消失。
\begin{yizhu}[易错点]
$\delta\SCS=0\Longrightarrow F=0$ 这一步默认了边界项消失。直接变分给出 $\delta\SCS=\frac{k}{2\pi}\int_M\Tr(\delta A\wedge F)$，再加边界项 $\frac{k}{4\pi}\int_{\partial M}\Tr(A\wedge\delta A)$；只有 $M$ 无边界，或限定 $\delta A$ 在 $\partial M$ 上为零，才能把后者丢掉而得到 $F^a_{\mu\nu}=0$。作者下文说全散度项“加上适当的边界条件”即消失，用的正是同一默认，但此处没有明说。
\end{yizhu}

% [Tutor Note: 易错点] (3.3) 的 2/3 与 Tr 归一化互相锁定：取基础表示 Tr(T^aT^b)=δ^{ab}/2，才有 (3.4)-(3.5) 的 ΔS=2πkw 与 w(g)∈Z。
\begin{yizhu}[易错点]
\eqref{eq:3.3} 中的 $\frac23$ 与 $\Tr$ 的归一化要分开看、配套用。$\frac23$ 这个系数保证作用量在规范变换下只差一个全导数，再加一项正比于 $\int\varepsilon^{\mu\nu\rho}\Tr(g^{-1}\partial_\mu g\,g^{-1}\partial_\nu g\,g^{-1}\partial_\rho g)$ 的量；$\Tr$ 的倍数则决定后者的数值：取基础表示 $\Tr(T^aT^b)=\frac12\delta^{ab}$ 时，\eqref{eq:3.5} 的 $w(g)$ 恰为整数，\eqref{eq:3.4} 的 $\frac13$ 与之配套给出 $\Delta\SCS=2\pi k\,w(g)$。若改取 $\Tr(T^aT^b)=\delta^{ab}$，$w(g)$ 整体翻倍，$1/24\pi^2$ 也得跟着改。作者随后正是拿“$w(g)$ 恰为整数”反过来定 $\Tr$ 的倍数。
\end{yizhu}

\paragraph{能级的量子化}
% 原文标题：Quantization of the level.
呼应 \S\ref{sec:2.2}，我们要问：陈--西蒙斯作用量 \eqref{eq:3.3} 在规范变换 \eqref{eq:3.1} 下是否不变？一次简短的分部积分练习给出
\begin{equation}\tag{3.4}\label{eq:3.4}
  \delta\SCS = \frac{k}{4\pi}\int_M \dd^3x\,\varepsilon^{\mu\nu\rho}\Bigl(\partial_\mu\Tr\bigl[(\partial_\nu g)(g^{-1}A_\rho)\bigr] + \frac{1}{3}\Tr\bigl[(g^{-1}\partial_\mu g)(g^{-1}\partial_\nu g)(g^{-1}\partial_\rho g)\bigr]\Bigr).
\end{equation}

% ==== 原文 p.15 ====
第一项是熟悉的全散度 \eqref{eq:2.10} 的乔装；与阿贝尔情形一样，只要加上适当的边界条件，它就消失。第二项是非阿贝尔理论的新东西，与一个称为 $g$ 的绕数的积分有关：
\begin{equation}\tag{3.5}\label{eq:3.5}
  w(g) \equiv \frac{1}{24\pi^2}\int_M \dd^3x\,\varepsilon^{\mu\nu\rho}\Tr\bigl[(g^{-1}\partial_\mu g)(g^{-1}\partial_\nu g)(g^{-1}\partial_\rho g)\bigr] = (\text{constant})\cdot n,
  \quad n\in\Z.
\end{equation}

% [Tutor Note: 易错点] (3.4) 的 1/3 与 (3.5) 的 1/24π² 必须配套：ΔS=(k/4π)(1/3)∫εTr=2πk·w，故 k∈Z 才是单值条件；任一系数改动都会破坏这一配对。
\begin{yizhu}[易错点]
\eqref{eq:3.4} 的 $\frac13$ 与 \eqref{eq:3.5} 的 $1/24\pi^2$ 是一对，不能分开改：规范变换后 $\Delta\SCS=\frac{k}{4\pi}\cdot\frac13\int_M\dd^3x\,\varepsilon^{\mu\nu\rho}\Tr\bigl[(g^{-1}\partial_\mu g)(g^{-1}\partial_\nu g)(g^{-1}\partial_\rho g)\bigr]=\frac{k}{12\pi}\int\cdots=2\pi k\,w(g)$。正因 $w(g)$ 取整数，振幅 $e^{i\SCS}$ 的单值性才归结为 $k\in\Z$。若把 $\frac13$ 改成别的数，$2\pi k\,w(g)$ 的整倍数结构就断了；\eqref{eq:3.3} 第二行那个差一倍的写法正是这种断裂的旁证。
\end{yizhu}

$w(g)$ 的量子化在瞬子研究中广为人知（证明见 \cite{coleman}）。因此我们固定 $\Tr$ 的归一化，使 $w(g)$ 恰为整数。用 $w(g)$ 表示，陈--西蒙斯作用量在规范变换下改变一个常数：$\SCS\to\SCS+2\pi k\,w(g)$。$w(g)=0$ 的变换保持 $\SCS$ 不变，$w(g)\neq0$ 的变换则不然。绕数非零的 $g(x)$ 称为大规范变换（large gauge transformation），它们非平凡地“缠绕” $G$。这类变换使经典陈--西蒙斯理论定义不良；但依 \S\ref{sec:2.2} 的论证，只要振幅 $e^{-\SECS}=e^{i\SCS}$ 单值，量子理论仍然自洽。这就要求 $2\pi k\,w(g)$ 为整数，从而 $k\in\Z$。

% [Tutor Note: 原文疑误] 「2πkw(g) 为整数」字面不成立：单值性要求 e^{2πikw}=1，即 k w(g)∈Z；结论 k∈Z 不变
\begin{yizhu}[原文疑误]
原文此句作 “This forces $2\pi k\,w(g)$ to be an integer”，字面读不通：$e^{i\SCS}$ 单值只要求 $e^{2\pi ik\,w(g)}=1$，即 $k\,w(g)\in\Z$，等价于说 $2\pi k\,w(g)$ 是 $2\pi$ 的整数倍，而不是它本身为整数（后者对一般的 $k$ 无解）。取 $w(g)=1$ 即得 $k\in\Z$，故结论不受影响。中译本照录原文措辞，读者按 $e^{2\pi ik\,w}=1$ 理解即可。
\end{yizhu}

\paragraph{规范变换的拓扑}
% 原文标题：Topology of gauge transformations.
$w(g)$ 的出现看似奇迹，其实是 $G$ 拓扑的推论。回忆：规范变换是映射 $g\colon M\to G$，$x\mapsto g(x)$。其中包含恒等映射 $x\mapsto 1_G$。人们可以问：是否所有规范变换都与恒等映射同伦（即可连续形变到恒等映射）？答案是否定的。例如取 $M=S^3$，映射 $g\colon S^3\to G$ 的集合由群 $\pi_3(G)$ 分类；博特（Bott）\cite{bott1956} 有一个著名结果：对每个紧致单李群 $G$，$\pi_3(G)=\Z$。

% [Tutor Note: 含糊点] π3(G)=Z 只引 Bott [12] 而不证：这是 Bott 用 Morse 理论研究紧李群拓扑的结论，生成元可取 S³≅SU(2)⊂G 的嵌入，其绕数正是 (3.5) 的 WZW 整数。
\begin{yizhu}[含糊点]
“$\pi_3(G)=\Z$” 出自博特 \cite{bott1956}，原文引而不证。这条定理来自 Bott 用莫尔斯理论研究紧李群的拓扑：对紧致单李群，$\pi_3$ 由一个生成元张成，该生成元可取 $S^3\cong\mathrm{SU}(2)\subset G$ 的嵌入。于是映射 $S^3\to G$ 的同伦类就是一个整数，而 \eqref{eq:3.5} 的积分正是这个整数的解析表达式（WZW 项的绕数）。注意定理要求 $G$ 紧致且单：阿贝尔群有 $\pi_3(\mathrm{U}(1))=0$，情形完全不同。
\end{yizhu}

因此每个映射 $g(x)$ 由整数 $w(g)=[g]\in\pi_3(G)=\Z$ 分类，这个整数称为绕数，它正是积分 \eqref{eq:3.5}。

\paragraph{威尔逊圈与纽结}
% 原文标题：Wilson loops and knots.
回忆：陈--西蒙斯理论是拓扑的，没有局部规范不变的可观测量，因为这类量会破坏广义协变性。它转而借助威尔逊圈计算拓扑不变量。非阿贝尔理论中的构造如下。固定 $G$ 的一个不可约表示 $R$，令 $\gamma$ 是 $M$ 中一条定向闭曲线，称为纽结。把联络 $A$ 看作 1-形式，沿 $\gamma$ 积分便得到 $\mathfrak{g}$ 的一个元素。该积分的（路径有序）指数给出 $G$ 的一个元素，它在共轭意义下确定：这个元素就是联络绕 $\gamma$ 的和乐。威尔逊圈 $W_R[\gamma]$ 则定义为该和乐在表示 $R$ 中的迹：
\begin{equation}\tag{3.6}\label{eq:3.6}
  W_R[\gamma] \equiv \Tr_R\Bigl[\mathcal{P}\exp\Bigl(\oint_\gamma \dd x^\mu A_\mu\Bigr)\Bigr].
\end{equation}

尽管表面上不然，这个定义从不需要在 $M$ 上选度规，因此 $W_R[\gamma]$ 是 $\gamma$ 的拓扑不变量。接着考虑 $M$ 中 $r$ 条定向且互不相交的闭曲线 $\gamma_1,\dots,\gamma_r$。它们的不交并 $L=\bigcup_{a=1}^r\gamma_a$ 称为链环。

% ==== 原文 p.16 ====
为每条曲线 $\gamma_a$ 选 $G$ 的一个不可约表示 $R_a$，并考虑下面的路径积分：
\begin{equation}\tag{3.7}\label{eq:3.7}
  \mathcal{Z}(M,L) \equiv \frac{1}{\mathcal{Z}}\int\mathcal{D}A\,e^{i\SCS[A]}\prod_{a=1}^{r}W_{R_a}[\gamma_a]
  = \Bigl\langle\prod_{a=1}^{r}W_{R_a}[\gamma_a]\Bigr\rangle.
\end{equation}

这里归一化因子 $\mathcal{Z}=\mathcal{Z}(M,\varnothing)$ 是在没有威尔逊圈时算出的。期望值 $\mathcal{Z}(M,L)$ 将是我们的主要研究对象：事实上，威滕（Witten）证明，当 $G=\mathrm{SU}(2)$、$M=S^3$ 时，$\mathcal{Z}(S^3,L)$ 恰是链环 $L$ 的琼斯多项式。

\paragraph{评注}
% 原文标题：Comments.
注意 $L$ 中任一圈 $\gamma$ 的定向，给出在相应表示 $R$ 下带荷的粒子绕该圈运动的方向。反转 $\gamma$ 的定向等价于对 $R$ 取共轭；因此同时作 $\gamma\to-\gamma$ 与 $R\to\bar R$，$\mathcal{Z}(M,L)$ 保持不变。此外，反转 $L$ 中所有圈的定向，等价于对所有表示取共轭。这个操作称为电荷共轭；它保持陈--西蒙斯作用量不变，也使 $\mathcal{Z}(M,L)$ 不变。

\subsection{弱耦合下的路径积分}\label{sec:3.2}

\paragraph{稳相}
% 原文标题：Stationary phase.
我们终于可以认真处理陈--西蒙斯路径积分了。先从弱耦合极限开始，它对应大 $k$。（确实，任意子相位与其他可观测量都依赖 $1/k$，$k$ 大时就变小。）没有威尔逊圈时，路径积分的被积函数快速振荡，因为 $\SCS$ 本身很大：
\begin{equation}\tag{3.8}\label{eq:3.8}
  \mathcal{Z} = \int\mathcal{D}A\,e^{i\SCS[A]}
  = \int\mathcal{D}A\,\exp\Bigl[\frac{ik}{4\pi}\int_M\Tr\Bigl(A\wedge\dd A + \frac{2}{3}A\wedge A\wedge A\Bigr)\Bigr].
\end{equation}

弱耦合极限下，$\mathcal{Z}$ 由稳相点的贡献主导。这类场位形 $A^{(\alpha)}$ 是经典运动方程的解，也就是平坦联络。为方便计算 $\SCS[A^{(\alpha)}]$，我们把 $k$ 提出来，定义纯拓扑的量 $I[A^{(\alpha)}]$，称为 $A^{(\alpha)}$ 的陈--西蒙斯不变量：
\begin{equation}\tag{3.9}\label{eq:3.9}
  I[A^{(\alpha)}] \equiv \frac{1}{4\pi}\int_M\Tr\Bigl(A^{(\alpha)}\wedge\dd A^{(\alpha)} + \frac{2}{3}A^{(\alpha)}\wedge A^{(\alpha)}\wedge A^{(\alpha)}\Bigr)
  \quad\Longrightarrow\quad \SCS[A^{(\alpha)}] = kI[A^{(\alpha)}].
\end{equation}

把每个鞍点贡献记为 $\mu[A^{(\alpha)}]$，我们对 $\mathcal{Z}$ 的首次估算写作
\begin{equation}\tag{3.10}\label{eq:3.10}
  \mathcal{Z} \approx \sum_\alpha \mu[A^{(\alpha)}] \approx \sum_\alpha e^{i\SCS[A^{(\alpha)}]} = \sum_\alpha e^{ikI[A^{(\alpha)}]}.
\end{equation}

这里指标 $\alpha$ 跑遍集合 $\mathcal{M}$，该集合称为平坦联络的模空间。暂时假定每个 $A^{(\alpha)}$ 孤立，且 $\mathcal{M}$ 有限；但一般情形并非如此。（我们将在 \S\ref{sec:4} 讨论：$\mathcal{M}$ 通常像个带若干奇点的流形。）这一练习的结果是：把所有平坦联络的贡献求和，$\mathcal{Z}$ 就算出 $M$ 的一个拓扑不变量。

% ==== 原文 p.17 ====
本节余下部分，我们考察每个 $A^{(\alpha)}$ 附近的小涨落，把上述计算做到一圈。

\paragraph{涨落}
% 原文标题：Fluctuations.
首先在“背景”平坦联络附近展开 $A=A^{(\alpha)}+a$，把这一分解代入 \eqref{eq:3.8}，再把路径积分的变量换成涨落场 $a$。作用量本身按 $a$ 的幂展开，最低的几项为
\begin{equation}\tag{3.11}\label{eq:3.11}
  \SCS[A] = kI[A^{(\alpha)}] + \frac{k}{4\pi}\int_M\Tr(a\wedge Da)
  = kI[A^{(\alpha)}] + \frac{k}{4\pi}\int_M\dd^3x\,\varepsilon^{\mu\nu\rho}\Tr(a_\mu D_\nu a_\rho).
\end{equation}

此后默认 $D_\mu=\partial_\mu+[A^{(\alpha)}_\mu,\cdot]$ 是关于背景平坦联络的协变导数。对每个这样的平坦联络 $A^{(\alpha)}$，路径积分 \eqref{eq:3.8} 中相应的鞍点贡献 $\mu[A^{(\alpha)}]$ 为
\begin{equation}\tag{3.12}\label{eq:3.12}
  \mu[A^{(\alpha)}] = e^{ikI[A^{(\alpha)}]}\int\mathcal{D}a\,\exp\Bigl[\frac{ik}{4\pi}\int_M\Tr(a\wedge Da)\Bigr]
  \equiv e^{ikI[A^{(\alpha)}]}\int\mathcal{D}a\,e^{iS[a]},
\end{equation}
这里已略去作用量的高阶项。这个路径积分是高斯型的；但要算它，必须先定规范，而这需要在 $M$ 上放一个度规。尽管如此，我们仍希望只用 $M$ 的纯拓扑不变量来算出 $\mathcal{Z}$。

\paragraph{规范固定}
% 原文标题：Gauge fixing.
我们取“洛伦兹”规范 $D_\mu a^\mu=0$。为实施它，我们采用法捷耶夫--波波夫规范固定手续。这一手续向路径积分引入一个辅助的拉格朗日乘子场 $\phi$，以及反对易却属玻色型的鬼场 $c$ 与 $\bar c$。（$\phi$、$\bar c$、$c$ 都是标量；但依维滕--巴尔--纳坦（Witten--Bar-Natan）\cite{barnatan1991} 的做法，我们把 $\phi$ 视为 3-形式。）为算出一圈贡献 $\mu[A^{(\alpha)}]$，我们给 \eqref{eq:3.12} 中的作用量 $S[a]$ 添上规范固定项与鬼场项，再对所有辅助场积分。先做第一步：把作用量改为
\begin{equation}\tag{3.13}\label{eq:3.13}
\begin{aligned}
  S[a] \longrightarrow S[a] + S_{\mathrm{gf}}[\phi,a] + iS_{\mathrm{gh}}[\bar c,c]
  &\equiv \int_M \dd^3x\,\Tr\Bigl(\frac{k}{4\pi}\varepsilon^{\mu\nu\rho}a_\mu D_\nu a_\rho + \phi D_\mu a^\mu + i\bar c D_\mu D^\mu c\Bigr)\\
  &= \int_M \Tr\Bigl(\frac{k}{4\pi}a\wedge Da + \phi\star D\star a + i\bar c D\star D c\Bigr).
\end{aligned}
\end{equation}

% [Tutor Note: 易错点] (3.13) 的鬼场项是 +i c̄D⋆Dc；经 (3.14) 外层再乘 i 后指数上是 −∫c̄Δc，故格拉斯曼积分给出 det(Δ) 在分子。漏掉任何一层 i 都会使 (3.17) 分子分母颠倒；c̄ 的上划线不可省。
\begin{yizhu}[易错点]
\eqref{eq:3.13} 括号内的鬼场项是 $+i\bar c\,D\star Dc$，注意是 $\bar c$ 而非 $c$，上划线不可省。到 \eqref{eq:3.14}，指数整体又乘一个 $i$，两层 $i$ 相乘得 $-1$，于是指数上的鬼场二次型变成 $-\int_M\Tr(\bar c D_\mu D^\mu c)$，正是 \eqref{eq:3.15} 左边那个减号。格拉斯曼高斯积分 $\int\mathcal{D}\bar c\,\mathcal{D}c\,e^{-\bar c\Delta c}=\det(\Delta)$，行列式落在分子；若漏掉 $S_{\mathrm{gh}}$ 前的 $i$ 或外层的 $i$，就会得到 $1/\det(\Delta)$，把 \eqref{eq:3.17} 的分子分母颠倒。
\end{yizhu}

于是每个平坦联络 $A^{(\alpha)}$ 的一圈贡献为
\begin{equation}\tag{3.14}\label{eq:3.14}
  \mu[A^{(\alpha)}] = e^{ikI[A^{(\alpha)}]}\int\mathcal{D}[a,\phi,\bar c,c]\,\exp\Bigl[i\int_M\Tr\Bigl(\frac{k}{4\pi}a\wedge Da + \phi\star D\star a + i\bar c D\star D c\Bigr)\Bigr].
\end{equation}

\paragraph{一圈行列式}
% 原文标题：One-loop determinants.
路径积分 \eqref{eq:3.14} 是高斯型的，可以拆成几块，分别用某些算子的行列式算出。鬼场只出现在动能项 $\bar c D_\mu D^\mu c$ 里，因此与作用量的其余部分解耦。

% ==== 原文 p.18 ====
对 $\bar c$ 与 $c$ 作格拉斯曼积分，得到动能算子的行列式：
\begin{equation}\tag{3.15}\label{eq:3.15}
  \int\mathcal{D}\bar c\,\mathcal{D}c\,\exp\Bigl[-\int_M\Tr\bigl(\bar c D_\mu D^\mu c\bigr)\Bigr] = \det(\Delta),
  \qquad \Delta\equiv D_\mu D^\mu = \star D\star D.
\end{equation}

路径积分的其余部分可以用一个称为 $L_-$ 的算子算出，其构造如下。先看算子 $L=\star D+D\star$。它有些像狄拉克算子\footnote{这是我们的第二个暗示——继 \S\ref{sec:2.4} 的评论之后——旋量的几何已经近在咫尺。}：其平方中含有拉普拉斯算子（还含若干曲率项）。在 3 维中，$L$ 自伴，并把偶次（相应地，奇次）形式映到偶次（相应地，奇次）形式；于是可以定义它到偶次形式与奇次形式上的限制 $L_+$ 与 $L_-$。接着把 $a$ 与 $\phi$ 合成为一个奇次形式场 $H=(a,\phi)$。可以证明：把 $a$ 与 $\phi$ 重新标度后，作用量 \eqref{eq:3.13} 的非鬼场部分恰是霍奇内积 $\frac12\langle H,L_-H\rangle$。它的路径积分因此是一个一圈行列式：
\begin{equation}\tag{3.16}\label{eq:3.16}
\begin{aligned}
  \mu_0[A^{(\alpha)}] &\equiv \int\mathcal{D}\phi\,\mathcal{D}a\,\exp\Bigl[i\int_M\Tr\Bigl(\frac{k}{4\pi}a\wedge Da + \phi\star D\star a\Bigr)\Bigr]\\
  &= \int\mathcal{D}H\,\exp\Bigl[\frac{i}{2}\int_M\Tr\bigl(H\wedge\star(L_-H)\bigr)\Bigr]
  = \frac{1}{\sqrt{\det(L_-)}}.
\end{aligned}
\end{equation}

% [Tutor Note: 易错点] (3.16) 前「L 把偶次映到偶次」是 3 维的 Z2 分次：⋆:Ω^k→Ω^{3−k}，故 ⋆D 会把 0-形式送到 2-形式；L+ 在 Ω0⊕Ω2、L− 在 Ω1⊕Ω3，H=(a,φ) 恰为奇次。
\begin{yizhu}[易错点]
“$L$ 把偶次（奇次）形式映到偶次（奇次）形式”应读作 $\Z_2$ 分次下的陈述，而不是保确切次数。3 维中 $\star$ 把 $k$-形式送到 $(3-k)$-形式：单独一个 $\star D$ 把 $0$-形式送到 $2$-形式，$D\star$ 把 $3$-形式送到 $1$-形式，两项相加后 $L$ 才保持分次。于是 $L_+$ 作用在 $\Omega^0\oplus\Omega^2$ 上，$L_-$ 作用在 $\Omega^1\oplus\Omega^3$ 上；$H=(a,\phi)$ 恰好是奇次的（$a$ 为 $1$-形式、$\phi$ 为 $3$-形式）。若误以为 $L_\pm$ 保确切次数，就会弄错 $\det(L_-)$ 所在的块与符号差。
\end{yizhu}

% [Tutor Note: 含糊点] (3.16) 的 it can be shown 跳过 ½⟨H,L−H⟩ 的拼装：用度规与 Hodge 星把 a∧Da 与 φ⋆D⋆a 的交叉项配成自伴算子 L−；细节见 Witten [1] 与 Bar-Natan–Witten [13]。
\begin{yizhu}[含糊点]
“可以证明”跳过了把两块拼成 $\frac12\langle H,L_-H\rangle$ 的步骤。实质是：先借 $M$ 的度规与霍奇星算子把 $a\wedge Da$ 与 $\phi\star D\star a$ 的交叉项配对，使非鬼场二次型成为一个自伴算子 $L_-$ 的霍奇内积；$\frac12$ 是把二次型写成标准霍奇内积时的归一化，缺了它 $\det(L_-)$ 会差一个整体因子。细节见维滕 \cite{witten1989}，亦可参见巴尔--纳坦与维滕 \cite{barnatan1991} 对 CS 微扰论中该算子的处理。
\end{yizhu}

把 \eqref{eq:3.14} 与 \eqref{eq:3.15}--\eqref{eq:3.16} 合起来，得到 $A^{(\alpha)}$ 的一圈贡献：
\begin{equation}\tag{3.17}\label{eq:3.17}
  \mu[A^{(\alpha)}] = e^{ikI[A^{(\alpha)}]}\mu_0[A^{(\alpha)}]\det(\Delta)
  = e^{ikI[A^{(\alpha)}]}\Biggl(\frac{\det(\Delta)}{\sqrt{\det(L_-)}}\Biggr).
\end{equation}

\paragraph{雷--辛格挠率}
% 原文标题：The Ray--Singer torsion.
必须承认，结果 \eqref{eq:3.17} 并不够显式。更令人担心的是，它还依赖 $M$ 的度规。尽管有这些顾虑，施瓦茨（Schwarz）\cite{schwarz1978} 仍有一个奇妙的结果：比值 $\det(\Delta)/\sqrt{\det(L_-)}$ 的绝对值是拓扑不变量，称为平坦联络 $A^{(\alpha)}$ 的雷--辛格挠率 $T_\alpha$。然而，这个比值的相位可能带来麻烦。先把该相位记为 $\theta_\alpha$，我们就可以把一圈配分函数 \eqref{eq:3.8} 写成鞍点贡献之和：
\begin{equation}\tag{3.18}\label{eq:3.18}
  \mu[A^{(\alpha)}] = e^{ikI[A^{(\alpha)}]}\bigl(T_\alpha e^{i\theta_\alpha}\bigr)
  \quad\Longrightarrow\quad
  \mathcal{Z} = \sum_\alpha e^{i(kI[A^{(\alpha)}] + \theta_\alpha)}T_\alpha.
\end{equation}

由于拉普拉斯算子 $\Delta$ 正定且自伴，实际上 $\det(\Delta)\in\R_+$。于是只剩研究 $\det(L_-)$ 的相位；我们将会看到，这一步相当棘手。

\paragraph{相位（其一）}
% 原文标题：The phase---I.
让我们换到 $L_-$ 的本征基，更仔细地考察路径积分 \eqref{eq:3.16}。本征函数记为 $\chi_j$，本征值记为 $\lambda_j$，即 $L_-\chi_j=\lambda_j\chi_j$。

% ==== 原文 p.19 ====
把上面引入的场 $H$ 展开为 $H=\sum_j h_j\chi_j$，得到
\begin{equation}\tag{3.19}\label{eq:3.19}
  \mu_0[A^{(\alpha)}] = \int\mathcal{D}H\,\exp\Bigl[\frac{i}{2}\int_M\Tr\bigl(H\wedge\star(L_-H)\bigr)\Bigr]
  = \prod_j\int_{-\infty}^{\infty}\frac{\dd h_j}{\sqrt{2\pi}}\,e^{i\lambda_jh_j^2},
\end{equation}
这里指数里的重复指标不求和。这次换元产生的积分剧烈振荡，说明 $\det(L_-)$ 的相位发散。不过它可以正规化：在一维情形，相应的积分是
\begin{equation}\tag{3.20}\label{eq:3.20}
  \int_{-\infty}^{\infty}\frac{\dd x}{\sqrt{2\pi}}\,e^{i\lambda x^2}
  = \lim_{\varepsilon\to0}\int_{-\infty}^{\infty}\frac{\dd x}{\sqrt{2\pi}}\,e^{i\lambda x}e^{-\varepsilon x^2}
  = \Bigl|\frac{1}{\sqrt{\lambda}}\Bigr|\exp\Bigl(\frac{i\pi}{4}\sgn(\lambda)\Bigr).
\end{equation}

% [Tutor Note: 原文疑误] (3.19)(3.20) 都漏了指数上的 ½：由 (3.16) 的 ½⟨H,L−H⟩ 应为 e^{(i/2)λh²}；补 ½ 后 (3.20) 右端 |1/√λ| 精确成立，不残留 1/√2。另 (3.20) 中段 e^{iλx} 漏了平方。
\begin{yizhu}[原文疑误]
两处照录的写法要合看。由 \eqref{eq:3.16}，高斯指数应为 $\frac i2\lambda_jh_j^2$，而 \eqref{eq:3.19} 与 \eqref{eq:3.20} 都写成不带 $\frac12$ 的 $i\lambda x^2$ 型。补上 $\frac12$ 后，$\int\dd x/\sqrt{2\pi}\,e^{i\lambda x^2/2}=|\lambda|^{-1/2}e^{i\pi\sgn(\lambda)/4}$ 与 \eqref{eq:3.20} 右端精确一致，并不残留 $1/\sqrt2$；所以此处并非“$\sqrt{2\pi}$ 测度被吸收”的问题，而是指数漏了 $\frac12$。若照字面按 $e^{i\lambda x^2}$ 算，右端应是 $|1/\sqrt{2\lambda}|$。另外 \eqref{eq:3.20} 中段的正规化因子漏写平方：$e^{i\lambda x}$ 应为 $e^{i\lambda x^2}$，否则极限不是 Fresnel 积分。中译本照录原文。
\end{yizhu}

对 \eqref{eq:3.19} 取乘积，得到 $\det(L_-)$ 的模与相位：
\begin{equation}\tag{3.21}\label{eq:3.21}
  \mu_0[A^{(\alpha)}] = \prod_j\Bigl|\frac{1}{\sqrt{\lambda_j}}\Bigr|\exp\Bigl(\frac{i\pi}{4}\sgn(\lambda_j)\Bigr)
  = \Bigl|\frac{1}{\sqrt{\det(L_-)}}\Bigr|\exp\Bigl(\frac{i\pi}{4}\sum_j\sgn(\lambda_j)\Bigr).
\end{equation}

于是行列式的相位取决于 $L_-$ 的符号差，即正本征值个数与负本征值个数之差。这个数定义不良，但引入 $\eta$ 不变量（eta invariant）\cite{aps1975,aps1976a,aps1976b} 就可以（又一次！）把它正规化：
\begin{equation}\tag{3.22}\label{eq:3.22}
  \mu_0[A^{(\alpha)}] = \frac{1}{\sqrt{\det(L_-)}}
  = \Bigl|\frac{1}{\sqrt{\det(L_-)}}\Bigr|e^{i\pi\eta[A^{(\alpha)}]/2},
  \qquad
  \eta[A^{(\alpha)}] \equiv \frac{1}{2}\lim_{s\to0}\sum_{\lambda\neq0}\frac{\sgn(\lambda)}{|\lambda|^s}.
\end{equation}

% [Tutor Note: 易错点] (3.22) 的 η 定义自带 ½，等于 APS 约定 η_APS=Σsign(λ)|λ|^{−s} 的一半，故 e^{iπη/2}=e^{iπη_APS/4} 才与 (3.21) 的 π/4 Σsign 相接。
\begin{yizhu}[易错点]
\eqref{eq:3.22} 的 $\eta$ 定义里已含 $\frac12$，故本节这个量等于 APS 约定的一半，即 $\eta=\frac12\eta_{\mathrm{APS}}$，其中 $\eta_{\text{APS}}(s)=\sum_{\lambda\neq0}\sgn(\lambda)|\lambda|^{-s}$ 见 \cite{aps1975}。正因如此，$e^{i\pi\eta/2}=e^{i\pi\eta_{\text{APS}}/4}$ 才与 \eqref{eq:3.21} 的 $\frac{\pi}{4}\sum_j\sgn(\lambda_j)$ 相接，\eqref{eq:3.23} 左边的 $\frac12$ 也按同一约定理解。若照搬 APS 文献的 $\eta$ 公式而丢掉这个 $\frac12$，相位会整体差一倍。
\end{yizhu}

\paragraph{相位（其二）}
% 原文标题：The phase---II.
为显式计算相位，阿蒂亚--帕托迪--辛格（Atiyah--Patodi--Singer）指标定理给出 $\eta[A^{(\alpha)}]$ 与规范丛其他拓扑数据的关系：
\begin{equation}\tag{3.23}\label{eq:3.23}
  \frac{1}{2}\bigl(\eta[A^{(\alpha)}] - \eta[0]\bigr) = \frac{c_2(G)}{2\pi}I[A^{(\alpha)}]
  \quad\Longleftrightarrow\quad
  \theta_\alpha = \frac{\pi}{2}\eta[A^{(\alpha)}] = \frac{c_2(G)}{2}I[A^{(\alpha)}] + \frac{\pi}{2}\eta[0].
\end{equation}

% [Tutor Note: 含糊点] (3.23) 引 APS 指标定理 [15–17] 而不证：取以 M 为边界的 4-维流形 X，对 X 上的 Dirac 算子用指标定理，边界修正即 η 项，体积分项用 ∫_X TrF∧F 化为 M 上的 I[A]。
\begin{yizhu}[含糊点]
\eqref{eq:3.23} 是阿蒂亚--帕托迪--辛格指标定理 \cite{aps1975,aps1976a,aps1976b} 的推论，原文只引不证。逻辑是：取一张以 $M$ 为边界的 4 维流形 $X$，对 $X$ 上的 Dirac 算子用指标定理；体积分项由 $\int_X\Tr F\wedge F$ 表出，边界修正正是 $\frac12(\eta[A^{(\alpha)}]-\eta[0])$ 这样的 $\eta$ 项。再用 $\int_X\Tr F\wedge F$ 与 $M$ 上陈--西蒙斯不变量 $I[A^{(\alpha)}]$ 的关系，即得 \eqref{eq:3.23}。式中 $\frac12$ 与 $c_2(G)$ 的归一化都须与 \eqref{eq:3.22} 的 $\eta$ 约定对齐。
\end{yizhu}

这里 $I[A^{(\alpha)}]$ 是 \eqref{eq:3.9} 定义的陈--西蒙斯不变量，$c_2(G)$ 是规范群的二次卡西米尔，神秘的 $\eta[0]$ 项是平凡联络的 $\eta$ 不变量。把 \eqref{eq:3.23} 与 \eqref{eq:3.18} 结合，得到一圈配分函数：
\begin{equation}\tag{3.24}\label{eq:3.24}
  \mathcal{Z} = e^{i\pi\eta[0]/2}\sum_\alpha\exp\Bigl[i\Bigl(k + \frac{c_2(G)}{2}\Bigr)I[A^{(\alpha)}]\Bigr]T_\alpha.
\end{equation}

% [Tutor Note: 物理图像] (3.24) 里 η 的贡献把能级从 k 平移到 k+c2(G)/2；这正是后文 q=e^{2πi/(k+N)} 与中心荷 c=kd/(k+N) 的起点。
\begin{yizhu}[物理图像]
\eqref{eq:3.24} 里 $\eta$ 的贡献除整体相位外，还把能级从 $k$ 平移到 $k+\frac{c_2(G)}{2}$。这一步不是技术细节：后文把 $k$ 换成 $k+N$（如 $\mathrm{SU}(N)$ 的 $q=e^{2\pi i/(k+N)}$）以及中心荷 $c=\frac{kd}{k+N}$，都以此平移为起点。读 \S\ref{sec:5} 时可回头对照。
\end{yizhu}

% =============================================================================
%  【片段汇报】s03a-framing.tex
%  覆盖：原文 p.13（"Apologia for geometry."）至 p.19 式 (3.24)；编号公式 (3.1)–(3.24) 共 24 个。
%  脚注：4、5（p.13）与 6（p.18）全译；页锚点 p.13–p.19 各一。
%  译注：13 条 = 易错点 7（(3.2)、(3.3) 归一化、(3.4)-(3.5)、(3.3) 变分边界、(3.13)、(3.16)、(3.22)）
%        + 含糊点 3（Bott [12]、"it can be shown" [1][13]、APS [15–17]）
%        + 原文疑误 2（(3.3) 第二行、(3.19)/(3.20) 漏 ½）+ 物理图像 1（(3.24) 能级平移）。
%
%  SPEC §10 第 6–10 条裁定（均已逐符号核对 p.14、p.18、p.19 页面图像）：
%   第 6 条 (3.4)：第一项确为全导数 ∂_μTr[(∂_νg)(g^{-1}A_ρ)]（注意是 ∂_νg，不是 g^{-1}∂_νg）；
%      第二项 1/3 与 (3.5) 的 1/24π² 配套 ⇒ ΔS=(k/12π)∫εTr=2πk·w，故 k∈Z，自洽。
%      但新发现 (3.3) 第二行差整体因子 2（见上「原文疑误」），与 (3.4) 的 1/3 矛盾。
%   第 7 条 (3.16)：L=⋆D+D⋆ 在 3 维保持 Z2 分次（⋆:Ω^k→Ω^{3−k} 翻转奇偶，D 也翻转，两次相消）；
%      L+ 作用于 Ω0⊕Ω2、L− 作用于 Ω1⊕Ω3，H=(a,φ) 为奇次。作者「even/odd degree」无误，已批注。
%   第 8 条 (3.20)：「1/√2 被每模 √(2π) 吸收」的解释不成立。真因是 (3.19)、(3.20) 漏了指数上的 ½
%      （由 (3.16) 的 ½⟨H,L−H⟩）；补 ½ 后 |1/√λ| 与 1/√det(L−) 精确成立，无 1/√2 残留；
%      照字面 e^{iλx²} 则右端应为 |1/√(2λ)|。另 (3.20) 中段 e^{iλx} 漏平方。已批注。
%   第 9 条 (3.22)：η 定义含 ½，即 η_本节=½η_APS；故 e^{iπη/2}=e^{iπη_APS/4} 与 (3.21) 的 π/4Σsign 相接，
%      与 (3.23) 的 ½ 也自洽。已核，属约定而非笔误（易错点已说明）。
%   第 10 条（p.20 η_g=η_g）：属 s03b 片段，本片段不涉及。
%
%  任务点名项的附加核对：
%   (3.2) [A_μ,A_ν] 系数为 1（对比常见 1/2 约定）：作者自洽，已批注。
%   (3.3) 2/3 因子与 Tr 归一化（基础表示 Tr(T^aT^b)=δ^{ab}/2）：配套成立，已批注。
%   (3.13) ics 项即 +i c̄D⋆Dc：符号正确；外层 i 再乘后成 −c̄Δc，故 det(Δ) 在分子，c̄ 的上划线不可省，已批注。
%   (3.20) 测度 1/√(2π)：见第 8 条裁定。
%   (3.22) η 定义含 1/2：见第 9 条裁定。
%   24 个 \tag 与原文编号、顺序逐一对应；p.13–p.19 页锚点齐全；\eqref 未硬编码编号。
%
%  新发现的原文笔误（已核页面图像；纯英文拼写/词序类只在汇报中记，正文照译不批注）：
%    p.13 "prevous section" → previous；p.13 "formulating nonabelian the CS theory" 词序应为 "the nonabelian CS theory"。
%    p.15 "may be be integrated" 重复 be；p.15 "2πkw(g) to be an integer" 严格应为「2π 的整数倍」
%         （按作者 §2 的简写习惯处理，未批注，特记于此）。
%    p.18 "its number of its positive eigenvalues" 重复 its。
%    p.19 "repreated" → repeated。
%    (3.3) 第二行整体差因子 2；(3.20) 中段漏平方（二者已在正文批注）。
%
%  存疑：无 [待确认]。p.14 "the trace of this matrix is proportional to Tr(ω)" 中的 ω 疑应为 ω_CS，
%        属记号小疵，未批注（不影响本段推导）。
% =============================================================================
